\documentclass[11pt]{article}
\usepackage[margin=1in]{geometry}
\usepackage{amsmath,amssymb,amsthm}
\usepackage{booktabs,longtable}
\usepackage{bm}

\newcommand{\R}{\mathbb{R}}
\newcommand{\x}{\mathbf{x}}
\newcommand{\bmu}{\boldsymbol{\mu}}
\newcommand{\argmin}{\mathop{\operatorname{arg\,min}}}
\newcommand{\clip}{\operatorname{clip}}
\newcommand{\med}{\operatorname{med}}
\newcommand{\Var}{\operatorname{Var}}
\newcommand{\st}{\operatorname{st}}
\newcommand{\ft}[1]{\texttt{\detokenize{#1}}}

\newtheorem*{prop}{Proposition}

\title{K-means formulation of the block-group clustering\\
\large (\texttt{standardize.py} + \texttt{cluster.py}, default settings)}
\date{}

\begin{document}
\maketitle

\section{The textbook objective}

Given observations $\x_1,\dots,\x_n\in\R^d$, k-means seeks
\begin{equation}
\argmin_{\mathbf S}\sum_{i=1}^{k}\sum_{\x\in S_i}\bigl\|\x-\bmu_i\bigr\|^2
=\argmin_{\mathbf S}\sum_{i=1}^{k}|S_i|\Var S_i,
\qquad
\bmu_i=\frac{1}{|S_i|}\sum_{\x\in S_i}\x .
\label{eq:wiki}
\end{equation}
The rest of this document defines exactly what $n$, $d$, $\x$, $\bmu_i$ and
$\|\cdot\|$ are in our pipeline.

%======================================================================
\section{Index sets}

\begin{itemize}
\item $\mathcal B$: all 2020 Census block groups (BGs) in the ACS 2019--23 universe,
  50 states + DC; $|\mathcal B| = 239{,}741$. Each $b\in\mathcal B$ has a 12-character GEOID;
  $\st(b)$ is its first 2 characters (state FIPS).
\item Row filter (the \emph{fit set}):
  \[
  \mathcal F=\{\,b\in\mathcal B:\ \mathrm{pop}_b\ge 50\ \text{and}\ \mathrm{households}_b>0\,\},
  \qquad |\mathcal F| = 236{,}829 .
  \]
\item Region $R$ (a set of state FIPS codes chosen with \texttt{--states}):
  \[
  \mathcal B_R=\{\,b\in\mathcal F:\ \st(b)\in R\,\},\qquad n=|\mathcal B_R| .
  \]
  For $R=\{09,34,36\}$ (CT, NJ, NY): $n=24{,}680$.
\item Features $\mathcal J$, $d=|\mathcal J|=42$ (Table~\ref{tab:features}).
\item Concepts $\mathcal G$, $|\mathcal G|=15$: each feature belongs to exactly one concept
  $g(j)\in\mathcal G$; $n_g=|\{j:g(j)=g\}|$ is the concept's column count.
\item $\mathcal E=\{\ft{ecef_x_km},\ft{ecef_y_km},\ft{ecef_z_km}\}\subset\mathcal J$ is the
  \emph{distance} concept ($n_{\text{distance}}=3$).
\item Pooled-scale blocks
  $\mathcal P=\{\mathcal E,K_{\text{race}},K_{\text{units}},K_{\text{lang}},K_{\text{occ}}\}$:
  disjoint sets of columns that share one scale (Steps 3--4). $K_{\text{race}}$ is the 7 race
  shares, $K_{\text{units}}$ the 6 home-type shares, $K_{\text{lang}}$ the 5 tract
  language shares (not \ft{lep_hh_share}), and $K_{\text{occ}}$ the 5 occupation shares. Every
  other feature is scaled on its own.
\end{itemize}

%======================================================================
\section{Constructing $\x_b$}

\subsection{Step 0: raw values}

$r_{b,j}$ is the raw value of feature $j$ for BG $b$ (possibly missing). The three distance
features are Earth-Centered Earth-Fixed coordinates of the BG's population-weighted center
(latitude $\varphi_b$, longitude $\lambda_b$; TIGER internal point as fallback) on the WGS84
ellipsoid at height $0$:
\begin{gather*}
a = 6378.137\ \text{km},\qquad f=\tfrac{1}{298.257223563},\qquad e^2=f(2-f),\qquad
N(\varphi)=\frac{a}{\sqrt{1-e^2\sin^2\varphi}},\\[4pt]
\mathbf p_b=\begin{pmatrix} r_{b,\mathrm{x}}\\ r_{b,\mathrm{y}}\\ r_{b,\mathrm{z}}\end{pmatrix}
=\begin{pmatrix}
N(\varphi_b)\cos\varphi_b\cos\lambda_b\\
N(\varphi_b)\cos\varphi_b\sin\lambda_b\\
N(\varphi_b)\,(1-e^2)\sin\varphi_b
\end{pmatrix}\ \text{km}.
\end{gather*}

\subsection{Step 1: per-feature transform}

With $\clip(v;\ell,u)=\min(\max(v,\ell),u)$,
\[
t_{b,j}=T_j(r_{b,j}),\qquad
T_j(v)=\begin{cases}
v & \text{none}\\[2pt]
\ln\!\bigl(1+\max(v,0)\bigr) & \text{log1p}\\[2pt]
\clip(v;\ell_j,u_j) & \text{winsor}\\[2pt]
\ln\!\bigl(1+\max(\clip(v;\ell_j,u_j),0)\bigr) & \text{winsor\_log1p}
\end{cases}
\]
The winsor bounds are empirical quantiles of the observed raw values over the national fit
set,
\[
\ell_j=\widehat Q_{\mathcal F}\!\bigl(r_{\cdot,j};\,p^{\mathrm{lo}}_j\bigr),\qquad
u_j=\widehat Q_{\mathcal F}\!\bigl(r_{\cdot,j};\,p^{\mathrm{hi}}_j\bigr),\qquad
(p^{\mathrm{lo}}_j,p^{\mathrm{hi}}_j)=\begin{cases}(0.020,\,0.995) & j=\ft{violent_crime_rate}\\
(0.005,\,0.995) & \text{otherwise.}\end{cases}
\]
Fitted values are listed in Table~\ref{tab:features}.

\subsection{Step 2: hierarchical median imputation}

Let $O_j\subseteq\mathcal F$ be the BGs where $t_{b,j}$ is observed. For a BG $b$ define nested
geographic neighbourhoods by GEOID prefix,
\[
A_L(b)=\{\,b'\in\mathcal F:\ \mathrm{GEOID}_{b'}[{:}\kappa_L]=\mathrm{GEOID}_b[{:}\kappa_L]\,\},
\qquad (\kappa_{\text{tract}},\kappa_{\text{county}},\kappa_{\text{state}},\kappa_{\text{nation}})=(11,5,2,0).
\]
For a scalar feature,
\[
\tilde t_{b,j}=\begin{cases}
t_{b,j} & b\in O_j\\
\med\{\,t_{b',j}: b'\in A_{L^*}(b)\cap O_j\,\} & b\notin O_j,
\end{cases}
\]
where $L^*$ is the first level in (tract, county, state, nation) whose donor set
$A_L(b)\cap O_j$ is non-empty.

\emph{Compositional blocks} $K\in\{$race (7), home type (6), tract language (5), occupation
(5)$\}$ are imputed all-or-nothing: if any share in $K$ is missing for $b$, the whole vector is
replaced. Donors are BGs whose block is complete, $O_K=\bigcap_{j\in K}O_j$, and
\[
\mathbf m_K(b)=\bigl(\med\{t_{b',j}:b'\in A_{L^*}(b)\cap O_K\}\bigr)_{j\in K},\qquad
\tilde t_{b,j}=\frac{m_{K,j}(b)}{\sum_{j'\in K}m_{K,j'}(b)}\quad(j\in K),
\]
so imputed shares still sum to 1. ($\ft{lep_hh_share}$ is in the language concept but not in
the tract-language block, so it is imputed as a scalar.)

\subsection{Step 3: national standardization (\texttt{features\_standardized.parquet})}

Population moments over the fit set ($\mathrm{ddof}=0$):
\[
m_j=\frac{1}{|\mathcal F|}\sum_{b\in\mathcal F}\tilde t_{b,j},\qquad
s_j=\Bigl(\frac{1}{|\mathcal F|}\sum_{b\in\mathcal F}(\tilde t_{b,j}-m_j)^2\Bigr)^{1/2}.
\]
Each column of a pooled-scale block $P\in\mathcal P$ is centered on its own mean, but all
columns of the block are divided by one shared scale: the root-mean of the block's column
variances.
\[
s_P=\Bigl(\frac{1}{|P|}\sum_{j\in P}s_j^2\Bigr)^{1/2},\qquad
s_{[j]}=\begin{cases}s_P & j\in P\in\mathcal P\\ s_j & \text{otherwise,}\end{cases}\qquad
z_{b,j}=\frac{\tilde t_{b,j}-m_j}{s_{[j]}} .
\]
Fitted values: $s_{\mathcal E}=862.33$ km, $s_{K_{\text{race}}}=0.1736$,
$s_{K_{\text{units}}}=0.1816$, $s_{K_{\text{lang}}}=0.1305$, $s_{K_{\text{occ}}}=0.1283$.
Each block keeps total variance $\sum_{j\in P}\Var(z_{\cdot,j})=|P|$, the same as $|P|$
separately z-scored features. Within the block, that variance is split among the columns in
proportion to each category's actual spread.

\subsection{Step 4: regional re-standardization (\texttt{--scale regional}, default)}

Restrict to $\mathcal B_R$ and standardize again, using regional moments:
\[
\bar z^R_j=\frac1n\sum_{b\in\mathcal B_R}z_{b,j},\qquad
\sigma^R_j=\Bigl(\frac1n\sum_{b\in\mathcal B_R}(z_{b,j}-\bar z^R_j)^2\Bigr)^{1/2},\qquad
\sigma^R_P=\Bigl(\frac{1}{|P|}\sum_{j\in P}(\sigma^R_j)^2\Bigr)^{1/2}\quad(P\in\mathcal P),
\]
\[
\check z_{b,j}=\begin{cases}
(z_{b,j}-\bar z^R_j)/\sigma^R_P & j\in P\in\mathcal P\\[2pt]
(z_{b,j}-\bar z^R_j)/\sigma^R_j & \text{otherwise, }\sigma^R_j>0\\[2pt]
z_{b,j}-\bar z^R_j\ (=0) & \text{otherwise, }\sigma^R_j=0 .
\end{cases}
\]

\begin{prop}[the national step cancels]
Because Step 3 is an affine map with positive slope, Step 4 equals standardizing the imputed,
transformed values directly within the region, using the same pooled blocks. Write
$\bar t^R_j$ and $\tau^R_j$ for the regional mean and population s.d.\ of $\tilde t_{\cdot,j}$,
and define the \emph{effective regional scale}
\[
\tau^R_{[j]}=\begin{cases}
\tau^R_P=\Bigl(\dfrac{1}{|P|}\sum_{j'\in P}(\tau^R_{j'})^2\Bigr)^{1/2} & j\in P\in\mathcal P\\[10pt]
\tau^R_j & \text{otherwise.}
\end{cases}
\]
Then
\[
\check z_{b,j}=\frac{\tilde t_{b,j}-\bar t^R_j}{\tau^R_{[j]}}\quad\forall j,
\qquad\text{in particular}\qquad
\begin{pmatrix}\check z_{b,\mathrm{x}}\\\check z_{b,\mathrm{y}}\\\check z_{b,\mathrm{z}}\end{pmatrix}
=\frac{\mathbf p_b-\bar{\mathbf p}^R}{\tau^R_{\mathcal E}},\quad
\tau^R_{\mathcal E}=\Bigl(\tfrac13\sum_{a\in\{\mathrm{x,y,z}\}}\Var_{\mathcal B_R}(p_{\cdot,a})\Bigr)^{1/2}.
\]
For $R=\{09,34,36\}$: $\tau^R_{\mathcal E}=94.88$ km (per-axis s.d.\ 126.6, 63.3, 83.5 km),
$\tau^R_{K_{\text{race}}}=0.1764$, $\tau^R_{K_{\text{units}}}=0.2219$,
$\tau^R_{K_{\text{lang}}}=0.1416$, $\tau^R_{K_{\text{occ}}}=0.1280$.
\end{prop}
\begin{proof}
For every $j$, $z_{b,j}=(\tilde t_{b,j}-m_j)/s_{[j]}$, so $\bar z^R_j=(\bar t^R_j-m_j)/s_{[j]}$ and
$\sigma^R_j=\tau^R_j/s_{[j]}$. Within a block $P$ the divisor $s_{[j]}=s_P$ is common, so
$\sigma^R_P=\tau^R_P/s_P$. Substituting gives the result.
\end{proof}

With \texttt{--scale national} instead, $\check z_{b,j}=z_{b,j}$ (no Step 4). Everything below
still holds with $(\bar t^R_j,\tau^R_{[j]})$ replaced by $(m_j,s_{[j]})$, except the
equal-contribution result of Section~\ref{sec:contrib}.

\subsection{Step 5: concept weights}

\[
w_j=\frac{c_{g(j)}}{\sqrt{n_{g(j)}}},\qquad c_g=1\ \ \forall g\in\mathcal G
\quad(\texttt{CONCEPT\_WEIGHTS}),\qquad
W=\operatorname{diag}(w_1,\dots,w_{42}).
\]
So $w_j=1$ for single-column concepts, $1/\sqrt2\approx0.7071$ for 2-column concepts,
$1/\sqrt3\approx0.5774$, $1/\sqrt5\approx0.4472$, $1/\sqrt6\approx0.4082$,
$1/\sqrt7\approx0.3780$.

\subsection{Step 6: the observation vector}

\[
\boxed{\;\x_b=W\check{\mathbf z}_b\in\R^{42},\qquad
x_{b,j}=w_j\,\check z_{b,j}=
\begin{cases}
\dfrac{c_{\text{distance}}}{\sqrt3}\cdot\dfrac{p_{b,a(j)}-\bar p^R_{a(j)}}{\tau^R_{\mathcal E}} & j\in\mathcal E\\[12pt]
\dfrac{c_{g(j)}}{\sqrt{n_{g(j)}}}\cdot\dfrac{\tilde t_{b,j}-\bar t^R_j}{\tau^R_{[j]}} & j\notin\mathcal E
\end{cases}\;}
\]
where $a(j)\in\{\mathrm x,\mathrm y,\mathrm z\}$ is the ECEF axis of feature $j$, and
$\tau^R_{[j]}$ is the block-pooled scale for share columns and the feature's own s.d.\
otherwise.

%======================================================================
\section{The objective with our definitions substituted}

The observations are $\{\x_b\}_{b\in\mathcal B_R}$, $n=|\mathcal B_R|$, $d=42$, $k$ from
\texttt{--k} (default 5). The partition is $\mathbf S=\{S_1,\dots,S_k\}$ of $\mathcal B_R$.

\paragraph{Centroids.} Since $\x_b$ is linear in $\check{\mathbf z}_b$ and in $\tilde{\mathbf t}_b$,
\[
\bmu_i=\frac1{|S_i|}\sum_{b\in S_i}\x_b=W\bar{\check{\mathbf z}}_{i},\qquad
\mu_{i,j}=\begin{cases}
\dfrac{c_{\text{distance}}}{\sqrt3}\cdot\dfrac{\bar p_{i,a(j)}-\bar p^R_{a(j)}}{\tau^R_{\mathcal E}} & j\in\mathcal E\\[12pt]
\dfrac{c_{g(j)}}{\sqrt{n_{g(j)}}}\cdot\dfrac{\bar t_{i,j}-\bar t^R_j}{\tau^R_{[j]}} & j\notin\mathcal E
\end{cases}
\]
with $\bar{\mathbf p}_i=\frac1{|S_i|}\sum_{b\in S_i}\mathbf p_b$ (the cluster's mean ECEF point)
and $\bar t_{i,j}=\frac1{|S_i|}\sum_{b\in S_i}\tilde t_{b,j}$.
(\texttt{centroids\_z.csv} reports $W^{-1}\bmu_i=\bar{\check{\mathbf z}}_i$.)

\paragraph{Objective.} Substituting into \eqref{eq:wiki}, the regional means cancel inside
$x_{b,j}-\mu_{i,j}$, leaving
\begin{align}
&\argmin_{\mathbf S}\sum_{i=1}^{k}\sum_{b\in S_i}\bigl\|\x_b-\bmu_i\bigr\|^2 \notag\\
&\quad=\argmin_{\mathbf S}\sum_{i=1}^{k}\sum_{b\in S_i}\sum_{g\in\mathcal G}\frac{c_g^2}{n_g}
\sum_{j:\,g(j)=g}\bigl(\check z_{b,j}-\bar{\check z}_{i,j}\bigr)^2 \notag\\
&\quad=\argmin_{\mathbf S}\sum_{i=1}^{k}\sum_{b\in S_i}\Biggl[
\frac{c_{\text{distance}}^2}{3\,(\tau^R_{\mathcal E})^2}\,\bigl\|\mathbf p_b-\bar{\mathbf p}_i\bigr\|_2^2
+\sum_{g\neq\text{distance}}\frac{c_g^2}{n_g}\sum_{j:\,g(j)=g}
\Bigl(\frac{\tilde t_{b,j}-\bar t_{i,j}}{\tau^R_{[j]}}\Bigr)^2\Biggr] \label{eq:obj}\\
&\quad=\argmin_{\mathbf S}\sum_{i=1}^{k}|S_i|\Var S_i,\qquad
\Var S_i=\frac1{|S_i|}\sum_{b\in S_i}\|\x_b-\bmu_i\|^2=\sum_{j=1}^{42}w_j^2\,\Var_{S_i}\!\bigl(\check z_{\cdot,j}\bigr). \notag
\end{align}
The first term is $\tfrac{1}{3(\tau^R_{\mathcal E})^2}$ times the sum of squared straight-line
(km) distances from each BG's ECEF point to its cluster's mean ECEF point.

\paragraph{Pairwise form.} Using
$\sum_{\x\in S}\|\x-\bmu\|^2=\frac{1}{2|S|}\sum_{\x,\mathbf y\in S}\|\x-\mathbf y\|^2$,
the same objective is
\[
\argmin_{\mathbf S}\sum_{i=1}^{k}\frac{1}{2|S_i|}\sum_{b\in S_i}\sum_{b'\in S_i}D^2(b,b'),
\]
with the pairwise distance $D$ defined next.

%======================================================================
\section{The distance function}

\subsection{General form}

\[
D(b,b')=\|\x_b-\x_{b'}\|_2,\qquad
D^2(b,b')=\sum_{j=1}^{42}w_j^2\bigl(\check z_{b,j}-\check z_{b',j}\bigr)^2
=\bigl(\tilde{\mathbf t}_b-\tilde{\mathbf t}_{b'}\bigr)^{\!\top}\Lambda\,\bigl(\tilde{\mathbf t}_b-\tilde{\mathbf t}_{b'}\bigr),
\]
\[
\Lambda=\operatorname{diag}\!\Bigl(\frac{w_j^2}{(\tau^R_{[j]})^2}\Bigr)_{j=1}^{42}.
\]
So $D$ is a \emph{diagonally weighted Euclidean distance} on the transformed, imputed features
$\tilde{\mathbf t}$. Each squared feature difference is divided by its concept's column count
and by a regional variance: the feature's own, or the block's pooled variance for ECEF
coordinates and share columns. Because whole blocks share one divisor, $\Lambda$ is not the
inverse of the diagonal of the covariance matrix, so $D$ is not a diagonal Mahalanobis
distance. There are no off-diagonal terms, so correlations between features are not adjusted
for.

\subsection{Decomposition by concept}

\[
D^2(b,b')=\sum_{g\in\mathcal G}\Delta^2_g(b,b'),\qquad
\Delta^2_g(b,b')=\frac{c_g^2}{n_g}\sum_{j:\,g(j)=g}\Bigl(\frac{\tilde t_{b,j}-\tilde t_{b',j}}{\tau^R_{[j]}}\Bigr)^2 .
\]
The terms below are written out with all $c_g=1$, using this notation:
\begin{itemize}
\item the superscript $R$ is dropped on $\tau$;
\item $\boldsymbol\pi^{K}_b$ is BG $b$'s share vector for block $K$, and $\pi$ on its own is a
  single share;
\item $\tau_K$ is the block's pooled regional scale;
\item $C(v)=\clip(v;\ell_j,u_j)$, using the bounds of the feature it is applied to.
\end{itemize}
\begin{align*}
\Delta^2_{\text{distance}}&=\frac{\|\mathbf p_b-\mathbf p_{b'}\|_2^2}{3\,\tau_{\mathcal E}^2}
 =\frac{\bigl(\text{chord}_{bb'}\ \text{km}\bigr)^2}{3\,\tau_{\mathcal E}^2}\\
\Delta^2_{\text{race}}&=\frac{\bigl\|\boldsymbol\pi^{\text{race}}_b-\boldsymbol\pi^{\text{race}}_{b'}\bigr\|_2^2}{7\,\tau_{K_{\text{race}}}^2}
 =\frac{1}{7\,\tau_{K_{\text{race}}}^2}\sum_{q\in\{\text{wht,blk,aian,asn,nhpi,oth,hisp}\}}
 \bigl(\pi^{\text{race}}_{b,q}-\pi^{\text{race}}_{b',q}\bigr)^2\\
\Delta^2_{\text{density}}&=\frac{1}{\tau_{\rho}^2}\Bigl(\ln\frac{1+\rho_b}{1+\rho_{b'}}\Bigr)^2
 \qquad(\rho=\text{people per km}^2)\\
\Delta^2_{\text{language}}&=\frac16\Biggl[
 \frac{\bigl\|\boldsymbol\pi^{\text{lang}}_b-\boldsymbol\pi^{\text{lang}}_{b'}\bigr\|_2^2}{\tau_{K_{\text{lang}}}^2}
 +\Bigl(\frac{\pi^{\text{LEP}}_b-\pi^{\text{LEP}}_{b'}}{\tau_{\text{LEP}}}\Bigr)^2\Biggr]
 \qquad(\boldsymbol\pi^{\text{lang}}\in\R^5:\ \text{eng, spa, ie, api, oth})\\
\Delta^2_{\text{age}}&=\Bigl(\frac{\text{age}_b-\text{age}_{b'}}{\tau_{\text{age}}}\Bigr)^2\\
\Delta^2_{\text{income}}&=\frac{1}{\tau_I^2}\Bigl(\ln\frac{1+I_b}{1+I_{b'}}\Bigr)^2
 \qquad(I=\text{median household income, \$})\\
\Delta^2_{\text{dependents}}&=\frac13\Biggl[
 \Bigl(\frac{\pi^{\text{kids}}_b-\pi^{\text{kids}}_{b'}}{\tau_{\text{kids}}}\Bigr)^2
 +\Bigl(\frac{C(h_b)-C(h_{b'})}{\tau_h}\Bigr)^2
 +\frac{1}{\tau_{\text{adr}}^2}\Bigl(\ln\frac{1+C(\text{adr}_b)}{1+C(\text{adr}_{b'})}\Bigr)^2\Biggr]\\
\Delta^2_{\text{home type}}&=\frac{\bigl\|\boldsymbol\pi^{\text{units}}_b-\boldsymbol\pi^{\text{units}}_{b'}\bigr\|_2^2}{6\,\tau_{K_{\text{units}}}^2}
 \qquad(\boldsymbol\pi^{\text{units}}\in\R^6:\ \text{1det, 1att, 2--4, 5--19, 20+, mobile})\\
\Delta^2_{\text{tenure}}&=\Bigl(\frac{\pi^{\text{rent}}_b-\pi^{\text{rent}}_{b'}}{\tau_{\text{rent}}}\Bigr)^2\\
\Delta^2_{\text{vehicles}}&=\frac12\Biggl[
 \Bigl(\frac{\pi^{\text{0veh}}_b-\pi^{\text{0veh}}_{b'}}{\tau_{\text{0veh}}}\Bigr)^2
 +\Bigl(\frac{C(v_b)-C(v_{b'})}{\tau_v}\Bigr)^2\Biggr]\\
\Delta^2_{\text{commute}}&=\frac{1}{\tau_m^2}\Bigl(\ln\frac{1+C(m_b)}{1+C(m_{b'})}\Bigr)^2
 \qquad(m=\text{median commute, min})\\
\Delta^2_{\text{education}}&=\frac12\Biggl[
 \Bigl(\frac{\text{SEDA}_b-\text{SEDA}_{b'}}{\tau_{\text{SEDA}}}\Bigr)^2
 +\Bigl(\frac{\text{HUD}_b-\text{HUD}_{b'}}{\tau_{\text{HUD}}}\Bigr)^2\Biggr]\\
\Delta^2_{\text{crime}}&=\frac12\Biggl[
 \frac{1}{\tau_{V}^2}\Bigl(\ln\frac{1+C(V_b)}{1+C(V_{b'})}\Bigr)^2
 +\frac{1}{\tau_{P}^2}\Bigl(\ln\frac{1+C(P_b)}{1+C(P_{b'})}\Bigr)^2\Biggr]
 \qquad(\text{per }100\text{k})\\
\Delta^2_{\text{property tax}}&=\Bigl(\frac{C(\text{tax}_b)-C(\text{tax}_{b'})}{\tau_{\text{tax}}}\Bigr)^2\\
\Delta^2_{\text{occupation}}&=\frac{\bigl\|\boldsymbol\pi^{\text{occ}}_b-\boldsymbol\pi^{\text{occ}}_{b'}\bigr\|_2^2}{5\,\tau_{K_{\text{occ}}}^2}
 \qquad(\boldsymbol\pi^{\text{occ}}\in\R^5:\ \text{mgmt, svc, sales, constr, prod})
\end{align*}
Every quantity here is the imputed value $\tilde t$ where the raw value was missing.

\subsection{How each kind of term behaves}

\begin{itemize}
\item \textbf{Geography.} $\Delta^2_{\text{distance}}$ is the squared straight-line distance
  through the Earth between the two population centers, in units of the region's pooled
  coordinate s.d. Because the three axes share one divisor, the term is rotation-invariant and
  isotropic: 10 km north counts the same as 10 km east. (Separate per-axis z-scores would
  multiply N--S and E--W distances by different factors.) On a sphere of radius $\rho_\oplus$,
  $\text{chord}=2\rho_\oplus\sin(\theta/2)$ for central angle $\theta$, which is within $0.1\%$
  of the great-circle distance up to $\sim700$ km and increases monotonically with it beyond
  that. In the tri-state run, two BGs 50 km apart get
  $\Delta^2_{\text{distance}}=50^2/(3\cdot94.88^2)\approx0.093$.
\item \textbf{Log-transformed features} (density, income, commute, age-dependency, crime).
  The term depends on the \emph{ratio} $\frac{1+v_b}{1+v_{b'}}$, not the difference, so
  doubling income adds the same distance at \$30k as at \$150k.
\item \textbf{Winsorized features.} Differences beyond $[\ell_j,u_j]$ are capped:
  $|C(v_b)-C(v_{b'})|\le u_j-\ell_j$.
\item \textbf{Compositional blocks} (race, home type, tract language, occupation) are compared
  by the plain Euclidean distance between the two share vectors, divided by one pooled scale
  per block. One percentage point therefore counts the same in every category of a block.
  Categories that vary more in the region still account for more of the block's variance;
  for example, the tri-state s.d.\ of the NH white share is 0.325 but the NHPI share's is
  0.005, so rare groups contribute little.
  Before the fix, each share was z-scored on its own. That divided rare shares by tiny
  s.d.s, so 1 pp of NHPI counted about $70\times$ as much as 1 pp of NH white, and a few BGs
  reached $|z|\approx50$--$67$.
  Single shares outside a block (renter, children, zero-vehicle, LEP) are still divided by
  their own s.d.
\item \textbf{Within-concept averaging.} Each concept's squared differences are
  \emph{averaged} over its $n_g$ columns, not summed. That stops a 7-column concept outweighing
  a 1-column concept.
\end{itemize}

\subsection{Every concept contributes equally on average}
\label{sec:contrib}

\begin{prop}
Under regional scaling, averaging over all ordered pairs in the region (including $b=b'$),
\[
\frac1{n^2}\sum_{b\in\mathcal B_R}\sum_{b'\in\mathcal B_R}\Delta^2_g(b,b')=2c_g^2
\quad\forall g,
\qquad
\frac1{n^2}\sum_{b,b'}D^2(b,b')=2\sum_{g}c_g^2=30 .
\]
\end{prop}
\begin{proof}
For any variable $y$, $\frac1{n^2}\sum_{b,b'}(y_b-y_{b'})^2=2\Var(y)$.

A feature outside every pooled block has regional variance $\Var(\check z_{\cdot,j})=1$. For a
pooled block $P$,
\[
\sum_{j\in P}\Var(\check z_{\cdot,j})=\frac{\sum_{j\in P}(\tau^R_j)^2}{(\tau^R_P)^2}=|P|.
\]
So each concept's columns have total regional variance $n_g$; the language concept, for
example, gets $5$ from $K_{\text{lang}}$ plus $1$ from LEP. Hence
$\frac{c_g^2}{n_g}\cdot n_g\cdot2=2c_g^2$.
\end{proof}
So with $c_g=1$, each of the 15 concepts accounts for $1/15$ of the average squared
distance. One distance concept's ``average'' contribution of 2 corresponds to a chord of
$\sqrt6\,\tau^R_{\mathcal E}\approx232$ km in the tri-state run. Raising $c_g$ to $\sqrt2$
doubles concept $g$'s share. This equality does \emph{not} hold under \texttt{--scale national}.

%======================================================================
\section{How the minimization is carried out}

\texttt{sklearn.cluster.KMeans(n\_clusters=k, n\_init=25, random\_state=seed)}
(scikit-learn 1.7; defaults \texttt{init="k-means++"}, \texttt{algorithm="lloyd"},
\texttt{max\_iter=300}, \texttt{tol=1e-4}) on the $n\times42$ matrix with rows $\x_b^\top$.

\begin{enumerate}
\item \textbf{Seeding (greedy k-means++).} Pick $\bmu_1^{(0)}$ uniformly from
  $\{\x_b\}$. Then, for $i=2,\dots,k$, draw $2+\lfloor\ln k\rfloor$ candidates with probability
  $\propto\min_{l<i}\|\x_b-\bmu^{(0)}_l\|^2$ and keep the candidate that most reduces the
  potential.
\item \textbf{Lloyd iterations} ($t=0,1,\dots$). Assign each BG to its nearest centroid:
  \[
  S_i^{(t)}=\Bigl\{b:\ \|\x_b-\bmu_i^{(t)}\|^2\le\|\x_b-\bmu_l^{(t)}\|^2\ \ \forall l\Bigr\}.
  \]
  Then update each centroid to its cluster mean:
  \[
  \bmu_i^{(t+1)}=\frac{1}{|S_i^{(t)}|}\sum_{b\in S_i^{(t)}}\x_b .
  \]
  Stop when the assignments stop changing, or when
  $\sum_i\|\bmu_i^{(t+1)}-\bmu_i^{(t)}\|^2\le10^{-4}\cdot\frac1{42}\sum_j\Var(x_{\cdot,j})$,
  or after 300 iterations. Under regional scaling
  $\frac1{42}\sum_j\Var(x_{\cdot,j})=\frac{1}{42}\sum_g c_g^2=15/42$.
\item \textbf{Restarts.} Repeat steps 1--2 25 times and keep the run with the lowest objective
  \eqref{eq:obj} (\texttt{inertia\_}). The result is a local minimum of
  \eqref{eq:wiki}, not a guaranteed global $\operatorname{arg\,min}$.
\item \textbf{Relabeling.} Renumber clusters $1,\dots,k$ so that
  $|S_1|\ge|S_2|\ge\dots\ge|S_k|$.
\item \textbf{Diagnostic.} The mean silhouette is computed on a uniform sample $\mathcal U$ of
  $\min(20000,n)$ BGs, using the same (unsquared) distance $D$:
  \[
  s(b)=\frac{\beta(b)-\alpha(b)}{\max\{\alpha(b),\beta(b)\}},\quad
  \alpha(b)=\operatorname*{mean}_{b'\in S_{i(b)}\cap\,\mathcal U\setminus\{b\}}D(b,b'),\quad
  \beta(b)=\min_{l\ne i(b)}\operatorname*{mean}_{b'\in S_l\cap\,\mathcal U}D(b,b').
  \]
\end{enumerate}

%======================================================================
\section{Feature table}

$w_j=c_{g(j)}/\sqrt{n_{g(j)}}$ with $c_g=1$. Winsor bounds are in raw units, fitted on
$\mathcal F$. $s_{[j]}$ is the national scale of $\tilde t$ from Step 3.

\begingroup\small
\setlength{\tabcolsep}{4pt}
\begin{longtable}{@{}llllrrrr@{}}
\caption{The $d=42$ features.}\label{tab:features}\\
\toprule
$j$ & feature & concept ($n_g$) & $T_j$ & $\ell_j$ & $u_j$ & $s_{[j]}$ & $w_j$\\
\midrule\endfirsthead
\toprule
$j$ & feature & concept ($n_g$) & $T_j$ & $\ell_j$ & $u_j$ & $s_{[j]}$ & $w_j$\\
\midrule\endhead
\bottomrule\endfoot
1 & \ft{ecef_x_km} & distance (3) & none & & & 862.33$^\dagger$ & 0.5774\\
2 & \ft{ecef_y_km} & distance (3) & none & & & 862.33$^\dagger$ & 0.5774\\
3 & \ft{ecef_z_km} & distance (3) & none & & & 862.33$^\dagger$ & 0.5774\\
4 & \ft{race_nh_white_share} & race (7) & none & & & 0.1736$^\dagger$ & 0.3780\\
5 & \ft{race_nh_black_share} & race (7) & none & & & 0.1736$^\dagger$ & 0.3780\\
6 & \ft{race_nh_aian_share} & race (7) & none & & & 0.1736$^\dagger$ & 0.3780\\
7 & \ft{race_nh_asian_share} & race (7) & none & & & 0.1736$^\dagger$ & 0.3780\\
8 & \ft{race_nh_nhpi_share} & race (7) & none & & & 0.1736$^\dagger$ & 0.3780\\
9 & \ft{race_nh_other_multi_share} & race (7) & none & & & 0.1736$^\dagger$ & 0.3780\\
10 & \ft{race_hispanic_share} & race (7) & none & & & 0.1736$^\dagger$ & 0.3780\\
11 & \ft{pop_density_per_sqkm} & density (1) & log1p & & & 2.0590 & 1.0000\\
12 & \ft{lang_tract_english_only_share} & language (6) & none & & & 0.1305$^\dagger$ & 0.4082\\
13 & \ft{lang_tract_spanish_share} & language (6) & none & & & 0.1305$^\dagger$ & 0.4082\\
14 & \ft{lang_tract_other_ie_share} & language (6) & none & & & 0.1305$^\dagger$ & 0.4082\\
15 & \ft{lang_tract_asian_pi_share} & language (6) & none & & & 0.1305$^\dagger$ & 0.4082\\
16 & \ft{lang_tract_other_share} & language (6) & none & & & 0.1305$^\dagger$ & 0.4082\\
17 & \ft{lep_hh_share} & language (6) & none & & & 0.0863 & 0.4082\\
18 & \ft{median_age} & age (1) & none & & & 10.270 & 1.0000\\
19 & \ft{median_hh_income} & income (1) & log1p & & & 0.5245 & 1.0000\\
20 & \ft{hh_with_children_share} & dependents (3) & none & & & 0.1543 & 0.5774\\
21 & \ft{avg_hh_size} & dependents (3) & winsor & 1.29 & 4.77 & 0.6117 & 0.5774\\
22 & \ft{age_dependency_ratio} & dependents (3) & winsor\_log1p & 0.0349 & 3.0321 & 0.1916 & 0.5774\\
23 & \ft{units_1_detached_share} & home type (6) & none & & & 0.1816$^\dagger$ & 0.4082\\
24 & \ft{units_1_attached_share} & home type (6) & none & & & 0.1816$^\dagger$ & 0.4082\\
25 & \ft{units_2_4_share} & home type (6) & none & & & 0.1816$^\dagger$ & 0.4082\\
26 & \ft{units_5_19_share} & home type (6) & none & & & 0.1816$^\dagger$ & 0.4082\\
27 & \ft{units_20plus_share} & home type (6) & none & & & 0.1816$^\dagger$ & 0.4082\\
28 & \ft{units_mobile_other_share} & home type (6) & none & & & 0.1816$^\dagger$ & 0.4082\\
29 & \ft{renter_share} & tenure (1) & none & & & 0.2733 & 1.0000\\
30 & \ft{zero_vehicle_hh_share} & vehicles (2) & none & & & 0.1359 & 0.7071\\
31 & \ft{mean_vehicles_per_hh} & vehicles (2) & winsor & 0.2137 & 3.0653 & 0.5207 & 0.7071\\
32 & \ft{median_commute_min} & commute (1) & winsor\_log1p & 7.8486 & 62.320 & 0.3429 & 1.0000\\
33 & \ft{school_quality_index} & education (2) & none & & & 0.3359 & 0.7071\\
34 & \ft{hud_school_proficiency_index} & education (2) & none & & & 28.692 & 0.7071\\
35 & \ft{violent_crime_rate} & crime (2) & winsor\_log1p & 16.103 & 1649.1 & 0.9922 & 0.7071\\
36 & \ft{property_crime_rate} & crime (2) & winsor\_log1p & 54.774 & 5155.7 & 0.8896 & 0.7071\\
37 & \ft{eff_property_tax_rate_tract} & property tax (1) & winsor & 0.001875 & 0.029171 & 0.005522 & 1.0000\\
38 & \ft{occ_mgmt_prof_share} & occupation (5) & none & & & 0.1283$^\dagger$ & 0.4472\\
39 & \ft{occ_service_share} & occupation (5) & none & & & 0.1283$^\dagger$ & 0.4472\\
40 & \ft{occ_sales_office_share} & occupation (5) & none & & & 0.1283$^\dagger$ & 0.4472\\
41 & \ft{occ_construction_maint_share} & occupation (5) & none & & & 0.1283$^\dagger$ & 0.4472\\
42 & \ft{occ_production_transp_share} & occupation (5) & none & & & 0.1283$^\dagger$ & 0.4472\\
\end{longtable}
\endgroup
\noindent$^\dagger$Pooled block scale $s_P$, shared by every column of the block: ECEF (km),
race, home-type units, tract language and occupation. Under regional scaling, the national
$s_{[j]}$ are replaced by the regional $\tau^R_{[j]}$. For CT+NJ+NY the pooled values are
$\tau^R_{\mathcal E}=94.88$ km, $\tau^R_{K_{\text{race}}}=0.1764$,
$\tau^R_{K_{\text{units}}}=0.2219$, $\tau^R_{K_{\text{lang}}}=0.1416$ and
$\tau^R_{K_{\text{occ}}}=0.1280$.

\end{document}
